% SEGMENT OLP-0523-S01
% Part: counterfactuals
% Chapter: minimal-change-semantics
% Section: introduction

\documentclass[../../../include/open-logic-section]{subfiles}

\begin{document}

\olfileid{cnt}{min}{int}

\olsection{அறிமுகம்}

``தீக்குச்சி உரசப்பட்டிருந்தால், அது பற்றியிருக்கும்'' போன்ற உண்மைக்கு
மாறான நிபந்தனைக் கூற்றுகளுக்கு ஸ்டால்நேக்கரும் லூயிஸும் விளக்கங்களை
முன்மொழிந்தனர். இத்தகைய வாக்கியங்களின் மெய்மை நிபந்தனைகளை எவ்வாறு சரியாகப்
புரிந்துகொள்வது என்பதற்கான முன்மொழிவுகளே அவர்களுடைய விளக்கங்கள். இரண்டின்
பின்னணிக் கருத்தும் இதுதான்: ஓர் உண்மைக்கு மாறான நிபந்தனைக் கூற்று
மெய்யா என்று மதிப்பிட, அதன் முன்னிலையை மெய்யாக்குவதற்குத் தேவையான அளவில்
மட்டுமே நடப்பு உலகிலிருந்து வேறுபடும் சாத்தியமான உலகங்களைக் கருத வேண்டும்.
அந்தச் சாத்தியமான உலகங்களில் பின்னிலை மெய்யாக இருந்தால், உண்மைக்கு மாறான
நிபந்தனையும் மெய்யாகும். எடுத்துக்காட்டாக, என் கையில் ஒரு தீக்குச்சியும்
தீப்பெட்டியும் இருப்பதாகக் கொள்க. நடப்பு உலகில் அவற்றை நான் வெறுமனே
பார்த்துக்கொண்டு, தீக்குச்சியை உரசினால் என்ன நடக்கும் என்று எண்ணுகிறேன்.
நான் தீக்குச்சியை உரசும் உலகுக்கும் நடப்பு உலகுக்கும் இடையிலான குறைந்தபட்ச
மாற்றம், நான் செயல்பட முடிவெடுத்துத் தீக்குச்சியை உரசுவதுதான். வேறெதுவும்
மாறாததால் அது குறைந்தபட்ச மாற்றம்: நான் அதே நேரத்தில் காற்றில்
குதிப்பதில்லை; தீக்குச்சியை உரசுவது என் தலைமுடியையும் பற்றவைப்பதில்லை;
என் விரல்களின் வலிமை திடீரென முற்றிலும் மறைவதில்லை; நீர்த் துப்பாக்கிப்
பதுங்குத் தாக்குதலில் அதே நேரத்தில் நான் நனைக்கப்படுவதில்லை; இப்படியே
பல. அந்த மாற்றுச் சாத்தியத்தில் தீக்குச்சி பற்றுகிறது. எனவே,
தீக்குச்சியை நான் உரசியிருந்தால் அது பற்றியிருக்கும் என்பது மெய்.

இந்த உள்ளுணர்வு விளக்கத்தை உண்மைக்கு மாறான நிபந்தனைகளின் தருக்கங்களுக்கான
முறைசார் பொருண்மையியலுடன் இணைக்கலாம். உண்மைக்கு மாறான நிபந்தனைக்கு லூயிஸ்
``$\boxright$'' என்ற குறியீட்டையும் ஸ்டால்நேக்கர்~``$>$'' என்ற
குறியீட்டையும் அறிமுகப்படுத்தினர். நாம்~$\cif$ ஐப் பயன்படுத்தி, அதை
கூற்றுத் தருக்கத்தில் ஓர் இரும இணைப்பியாகச் சேர்ப்போம். எனவே, $!A \lif
!B$ வடிவிலுள்ள !!{formula}s உடன் $!A \cif !B$ வடிவிலுள்ள !!{formula}s
உம் நமக்கு உண்டு. வாய்ப்புநிலைத் தருக்கத்திற்கான தொடர்புசார்
பொருண்மையியலைப் போலவே, முறைசார் பொருண்மையியலும் உலகங்களில் !!{formula}s
மதிப்பிடப்படும் மாதிரிகளை அடிப்படையாகக் கொண்டது. $\mSat{M}{!A \cif
!B}[w]$ ஐ வரையறுக்கும் நிறைவுறுத்தல் நிபந்தனை, சில (வேறு)
உலகங்கள்~$w'$ க்கான $\mSat{M}{!A}[w']$, $\mSat{M}{!B}[w']$ ஆகியவற்றைக்
கொண்டு தரப்படுகிறது. எந்த~$w'$? உள்ளுணர்வுப்படி, $\mSat{M}{!A}[w']$
ஆகும் உலகங்களில்~$w$ க்கு மிக நெருங்கிய ஒன்று அல்லது பல. எனவே,
``நெருக்கம்'' குறித்த தொடர்பு ஒன்றையும் மாதிரியில் சேர்க்க வேண்டும்.

உண்மைக்கு மாறான சூழ்நிலைகளைப் படமாகக் காட்டுவதற்கான விளக்கமிக்க ஒரு
முறையை லூயிஸ் அறிமுகப்படுத்தினார். ஒவ்வொரு சாத்தியமான உலகமும் மற்ற
உலகங்களைக் கொண்ட ஒன்றினுள் ஒன்று அடங்கிய கோளங்களின் கணமொன்றின் மையத்தில்
உள்ளது---இக்கோளங்களை நாம் ஒரே மையமுள்ள வட்டங்களாக வரைகிறோம். இரு
கோளங்களுக்கு இடையிலுள்ள உலகங்கள் ஒன்றுக்கொன்று மைய உலகுக்கு நிகரான
அளவு நெருக்கமானவை; உட்புறக் கோளத்தில் அடங்கியவை இன்னும் நெருக்கமானவை;
சுற்றியுள்ள கோளத்தில் இருப்பவை இன்னும் தொலைவானவை.
\begin{center}
\begin{tikzpicture}[scale=.7]\small
  \spheresystem{5}
  \draw (0,0) node {$w$};
  \propositionintersect{0}{4}{40}{3.5}
  {\tikzset{proposition/.append={smooth,tension=1.4}}
  \proposition[shift={(1.6,-1)}]{90}{1}{30}{4}}
  \path (1.5,3) node[below] {$\formula{B}$};
  \path (3,0) node {$\formula{A}$};
\end{tikzpicture}
\end{center}
மைய உலகம்~$w$ ஐச் சுற்றியுள்ள மிகச் சிறிய கோளத்தில் (சாம்பல் நிறப்
பகுதியில்) அமைந்து~$!A$ ஐ நிறைவுசெய்யும் உலகங்கள்~$w'$ தான் மிக
நெருங்கிய $!A$-உலகங்கள். உள்ளுணர்வுப்படி, மிக நெருங்கிய எல்லா
$!A$-உலகங்களிலும் $!B$ மெய்யாக இருந்தால்~$w$ இல் $!A \cif !B$
நிறைவுசெய்யப்படுகிறது.

\end{document}
